\documentclass[11pt]{article}
\usepackage{amsmath, amssymb}
\usepackage{geometry}
\geometry{margin=2.5cm}

\newcommand{\dd}{\mathrm{d}}
\newcommand{\Tr}{\mathrm{Tr}}
\newcommand{\calN}{\mathcal{N}}
\newcommand{\calO}{\mathcal{O}}

\title{Preconditioned Crank--Nicolson sampling for spinDMFT}
\date{}

\begin{document}
\maketitle

\section{Setting}

The spin-DMF\!T self-consistency equations require the thermal spin correlations
\begin{equation}
    g^{\alpha\beta}(\tau)
    = \frac{\displaystyle\int \dd V\, p(V)\,
             \Tr\!\bigl[U(\beta{-}\tau,\beta;V)\,S^\alpha\,U(\tau,0;V)\,S^\beta\bigr]}
           {\displaystyle\int \dd V\, p(V)\, Z(V)},
    \label{eq:g}
\end{equation}
where $V = \{V^\alpha(\tau_n)\}$ is the mean-field noise trajectory living in the
diagonal Matsubara frequency basis,
$p(V) = \prod_{n,\alpha}\calN(0,\lambda_{n,\alpha})$ is a Gaussian prior with
eigenvalues $\lambda_{n,\alpha}$ determined self-consistently from the covariance
matrix, $U(\tau_2,\tau_1;V)$ is the imaginary-time propagator generated by the
single-spin Hamiltonian $H(\tau;V) = V(\tau)\cdot S$, and
$Z(V) = \Tr[U(\beta,0;V)]$ is the per-trajectory partition function.

Defining the target measure
\begin{equation}
    \pi(\dd V) \propto p(V)\,Z(V)\,\dd V,
    \label{eq:target}
\end{equation}
equation~\eqref{eq:g} can be written as an expectation under $\pi$:
\begin{equation}
    g^{\alpha\beta}(\tau)
    = \mathbb{E}_\pi\!\left[
        \frac{\Tr[U(\beta{-}\tau,\beta;V)\,S^\alpha\,U(\tau,0;V)\,S^\beta]}{Z(V)}
      \right].
\end{equation}
An unbiased Monte Carlo estimator therefore accumulates the per-sample observable
$\Tr[\ldots]/Z(V)$ under draws $V^{(i)}\sim\pi$.

\section{Difficulties with random-walk Metropolis--Hastings}

The simplest Markov chain targeting $\pi$ uses Gaussian random-walk proposals
$V' = V + \varepsilon\,\xi$, $\xi\sim p(V)$, with Metropolis--Hastings
acceptance probability
\begin{equation}
    a(V\to V')
    = \min\!\left(1,\frac{p(V')Z(V')}{p(V)Z(V)}\right).
\end{equation}
In the diagonal Matsubara frequency basis the prior factorises, so the proposal
noise is drawn component-wise as $\xi_{n,\alpha}\sim\calN(0,\lambda_{n,\alpha})$.
The step size $\varepsilon$ is tuned for an acceptance rate of roughly $25\%$.

The fundamental limitation of random-walk M--H is that its autocorrelation time
scales as $\tau_{\mathrm{int}} = \calO(d)$, where $d$ is the dimension of $V$.
Here $d = 3N_\tau$ (three spin components times the number of time steps).
For $N_\tau = 50$ this gives $d = 150$, so only $N/\tau_{\mathrm{int}} \approx N/150$
effectively independent samples are obtained from $N$ M--H steps.
Because the importance-sampling estimator produces \emph{independent} draws from
$p(V)$, random-walk M--H offers no practical advantage at moderate $N_\tau$.

\section{Preconditioned Crank--Nicolson (pCN)}

The pCN algorithm~\cite{Cotter2013} is designed precisely for targets of the form
$\pi \propto p(V)\,\Phi(V)$ where $p$ is Gaussian and $\Phi(V) = Z(V) > 0$.
Its proposal reads
\begin{equation}
    \boxed{V' = \sqrt{1-\beta^2}\,V + \beta\,\xi, \qquad \xi \sim p(V),}
    \label{eq:pCN}
\end{equation}
with $\beta\in(0,1]$ a tunable step-size parameter.
The acceptance probability simplifies to
\begin{equation}
    a(V\to V') = \min\!\left(1,\frac{Z(V')}{Z(V)}\right),
    \label{eq:accept}
\end{equation}
because the Gaussian factors in $p(V')$ and $p(V)$ cancel exactly against the
proposal densities.  The key properties are:

\begin{enumerate}
    \item \textbf{Dimension-independent mixing.}  For targets of the form
    $\pi \propto p \cdot \Phi$ the acceptance rate and autocorrelation time of
    pCN are independent of $d$, in contrast to the $\calO(d)$ scaling of
    random-walk M--H.  This makes pCN the preferred algorithm in high-dimensional
    Gaussian-prior settings.

    \item \textbf{No prior-ratio computation needed.}  In the random-walk variant
    the term $p(V')/p(V)$ requires evaluating $\log p$, which involves summing
    over all $d$ components.  In pCN this ratio is absorbed into the proposal by
    construction, so only $Z(V')$ and $Z(V)$ need to be evaluated.  Since a
    $Z$-evaluation requires a full imaginary-time propagation (the bottleneck of
    the simulation), pCN and random-walk M--H cost the same per step while pCN
    mixes faster.

    \item \textbf{Single tuning parameter.}  The scalar $\beta$ controls the
    proposal step size.  $\beta = 1$ recovers independence M--H (draw $V'$
    fresh from the prior), $\beta \to 0$ gives infinitesimal moves.  Optimal
    performance is typically achieved at $\beta$ corresponding to an acceptance
    rate of $20$--$40\%$.
\end{enumerate}

\section{Implementation in the diagonal Matsubara basis}

Working in the basis that diagonalises the per-block covariance matrix (with
eigenvalues $\lambda_{n,\alpha}$), the prior is a product of independent
Gaussians.  In this basis the pCN proposal~\eqref{eq:pCN} acts component-wise:
\begin{equation}
    V'_{n,\alpha}
    = \sqrt{1-\beta^2}\,V_{n,\alpha}
      + \beta\,\xi_{n,\alpha},
    \qquad
    \xi_{n,\alpha} \sim \calN(0,\lambda_{n,\alpha}).
\end{equation}
After proposing $V'$ in the diagonal basis one applies the orthogonal
back-rotation and the inverse Fourier transform to obtain the time-domain field
$V'(\tau_n)$, which is then passed to the propagator routine.  The chain state
and the acceptance step~\eqref{eq:accept} are otherwise identical to the
random-walk variant already implemented in \texttt{Algorithm\_MH/}.

The main loop changes relative to the random-walk implementation are:
\begin{enumerate}
    \item Replace the random-walk proposal
    $V' = V + \varepsilon\,\xi$ with $V' = \sqrt{1-\beta^2}\,V + \beta\,\xi$.
    \item Remove the $\log p(V')/p(V)$ term from the log-acceptance ratio;
    only $\log Z(V') - \log Z(V)$ remains.
    \item Rename the CLI parameter \texttt{--mhStepSize} to \texttt{--pcnBeta}
    (or repurpose it) with a default of, e.g., $\beta = 0.5$.
\end{enumerate}

\section{Expected performance gain}

For the random-walk variant with $d = 150$ and $N = 80\,000$ nominal samples
the estimated autocorrelation time is $\tau_{\mathrm{int}} \approx d / (2.38^2
\times 0.234) \approx 115$ steps, giving an effective sample size
$\mathrm{ESS} \approx 700$.  Under pCN with the same $N$ and a comparable
acceptance rate the autocorrelation time is $\calO(1)$ in $d$, so
$\mathrm{ESS} \approx N$, representing a gain of roughly two orders of magnitude
at $N_\tau = 50$.  The gain grows linearly with $N_\tau$, making pCN
increasingly advantageous at finer imaginary-time discretisations.

\begin{thebibliography}{1}
\bibitem{Cotter2013}
S.~L.~Cotter, G.~O.~Roberts, A.~M.~Stuart, D.~White,
\textit{MCMC methods for functions: modifying old algorithms to make them faster},
Statistical Science \textbf{28}, 424 (2013).
\end{thebibliography}

\end{document}
